%
% einleitung.tex -- Beispiel-File für die Einleitung
%
% (c) 2020 Prof Dr Andreas Müller, Hochschule Rapperswil
%
% !TEX root = ../../buch.tex
% !TEX encoding = UTF-8
%
\section{Motivation (Noel) \label{laplace:section:teil0}}
\kopfrechts{Teil 0}

%\subsection{Laplacetransformation}
Die Laplace-Transformation ist eine Integral-Transformation von der Form
\begin{equation}
F(s) = \mathcal{L}\{f(t)\} = \int_{0}^{\infty} f(t) \, e^{-st} \, \mathrm{d}t ,
%\label{Laplace:teil0:Laplacetransformation}
\end{equation}
wobei $f(t)$ die zu transformierende Zeitfunktion ist ($f$ ist meistens Real) und $s$ ein Parameter der Transformation aus dem Komplexen Zahlenraum.
Der Real- und Imaginärteil von $s$ wird meistens mit $\delta$ und $\omega$ abgekürzt.
Somit repräsentiert mit $e^{-st} = e^{\delta t} \cdot e^{i \omega t}$ im Integral $\delta$ eine Abklingkonstante und $\omega$ eine Kreisfrequenz mit der auf Orthogonalität mit $f$ geprüft wird.
Des weiteren bildet $F(s)$ die Laplace-Transformierte.


Bei der Laplace-Transformierten $F(s)$ spricht man von einer Funktionen im Bildbereich die mit einer Funktion $f(t)$ in Korrespondenz steht.
Dabei spricht man $f(t)$ Korrespondiert mit $F(t)$.
Diese Korrespondenz wird mit einem "Hantelsymbol" geschrieben so dass $f(t) \laplacecorr F(s)$ oder $F(s) \invlaplacecorr f(t)$ Symbol ausgedrückt.

Als Beispiel wird die Funktion $f(t) = e^{\alpha t}$ transformiert.
$\alpha$ kann dabei beliebig Komplex sein, solange das Integral konvergiert.
Dies ist garantiert mit $\mathrm{Re}\{\alpha\} < \delta$.
Eingesetzt ergibt das
\begin{equation}
	\begin{aligned}
		\mathcal{L}\{f(t)\} &= \int_{0}^{\infty} e^{\alpha t} \, e^{-st} \, \mathrm{d}t \\
		&= \frac{1}{s - \alpha}.
	\end{aligned}
	\label{Laplace:teil0:exp beispiel}
\end{equation} 
Somit ist auch klar, dass $\mathcal{L}^{-1}\{\frac{1}{s - \alpha}\} = e^{\alpha t}$ ergibt.
Wir haben also die erste Korrespondenz,
\begin{equation}
	e^{\alpha t} \laplacecorr \frac{1}{s - \alpha},
\end{equation}
aufgestellt.
Kern dieses Kapitels ist es, solche Korrespondenzen, durch die Laplace-Umkehrformel, auch genannt Bromwich-Integral, durch spezielle Konturen aufzustellen.

Breite Anwendung findet die Laplace-Transformation vor allem dadurch, dass sie Ableitungen in algebraische Ausdrücke umwandeln kann.
Dadurch können sonst schwierige Differentialgleichungen vereinfacht oder direkt gelöst werden.
Dies kann einfach gezeigt werden indem die nach der Zeit abgeleitete Funktion $f'(t)$ eingesetzt und anschliessend durch Partielle Integration integriert wird.
Dadurch wird aus
\begin{equation}
	\mathcal{L}\{f'(t)\} = \int_{0}^{\infty} f'(t) \, e^{-st} \, \mathrm{d}t,
	%\label{Laplace:teil0:integral ableitung laplacetransformation}
\end{equation}
wobei $f(t)$ die Stammfunktion von $f'(t)$ $f$ ist,
\begin{equation}
	\mathcal{L}\{f'(t)\} = \left. e^{-st} f(t) \, \right|_{t=0}^{t=\infty} - 
	\int_{0}^{\infty} f(t) \, (-s \, e^{-st}) \, dt.
	%\label{Laplace:teil0:partielle Integration}
\end{equation}
Mit der Mit der Annahme, dass $\mathrm{Re}\{e^{-st}\} = e^{-\delta}$ stärker abschwächt, als $f(t)$ wächst, wird der erste Term gegen $0 - f(0^+)$ konvergieren. % TODO: konvergenz mit 0^+ kontrollieren
Der Zweite Term ist wiederum die normale Laplace-Transformation mit dem Vorfaktor $s$. 
Dadurch ergibt sich
\begin{equation}
	\begin{aligned}
		\mathcal{L}\{f'(t)\} &= s \mathcal{L}\{ f(t) \} - f(0^{+}) \\
		&= sF(s) - f(0^{+}).
	\end{aligned}
	\label{Laplace:teil0:ableitung laplacetransformation}
\end{equation}
Somit ergibt sich als erste Korrespondenz 
\begin{equation}
	f'(t) \laplacecorr sF(s) - f(0^+).
	%\label{Laplace:teil0:ableitung korrespondenz}
\end{equation}
Die Ableitung Korrespondiert nun, sofern $f(0^+)$ bekannt ist, mit einem algebraischen Ausdruck. 
Ableitungskorrespondenzen höheren Grades können einfach durch erneute  Substitution von $ f(t)$ mit $f'(t)$ in Gleichung \ref{Laplace:teil0:ableitung laplacetransformation} gefunden werden.
Dabei ergibt sich für die n-te Ableitung
\begin{equation}
	f^{(n)}(t) \laplacecorr s^{n}F(s) - s^{n-1}f(0^+) - s^{n-2}f'(0^+) - \dotsc - f^{(n-1)}(0^+).
\end{equation}

An 

